% Working-draft insertion, 2026-09-12: reviewed against the local Chapter 4.
% This file is an insertion source, not an independently compiled document.
\subsection{Working alternative: edge quadrature of the nonlinear physical flux}
\label{section:ALE_RD_Gauss_physical_flux}

\red{\textbf{Mathematical working note.} The following distinguishes three
spatial residual constructions. The quadrature identities are established
on an affine triangle with fixed connectivity during one residual evaluation;
their use in the complete moving-mesh method remains subject to tests of
accuracy, stability and topology changes. None is selected here solely
because its expression is simpler in an element co-moving frame.}

The $P^1$ assumption specifies
$\mathbf U_h=\sum_i\psi_i\mathbf U_i$. It does not require the nonlinear
Euler flux $\boldsymbol{\mathcal F}(\mathbf U_h)$ to be a $P^1$ function.
In particular,
\begin{equation}
 \boldsymbol{\mathcal F}\left(\sum_i\psi_i\mathbf U_i\right)
 \ne \sum_i\psi_i\boldsymbol{\mathcal F}(\mathbf U_i)
 \label{eq:ALE_RD_Gauss_noncommutation}
\end{equation}
in general. Equation~\eqref{eq:ALE_RD_P1U_contour_sum} exactly integrates
the specified nodal-flux interpolant on each straight edge. A different
choice is to approximate the integral of the nonlinear flux generated by
the same conservative $P^1$ state. Increasing the number of quadrature
points while retaining the same linear flux interpolant would not change
the old result.

Let edge $e=(a,b)$ have length $\ell_e$, outward unit normal $\mathbf n_e$,
and parameter $s\in[0,1]$. The edge state is
\begin{equation}
 \mathbf U_e(s)=(1-s)\mathbf U_a+s\mathbf U_b.
\end{equation}
The old and new spatial residuals are respectively
\begin{align}
 \widetilde\Phi_{\mathrm{nodal}}^T
 &=\sum_e\ell_e\left[
 \frac{\boldsymbol{\mathcal F}(\mathbf U_a)+
       \boldsymbol{\mathcal F}(\mathbf U_b)}{2}\cdot\mathbf n_e
 -(\overline{\boldsymbol\sigma}_T\cdot\mathbf n_e)
        \frac{\mathbf U_a+\mathbf U_b}{2}\right],
 \label{eq:ALE_RD_Gauss_old_edge}\\
 \widetilde\Phi_{G_q}^T
 &=\sum_e\ell_e\left[
 \sum_{r=1}^q w_r\boldsymbol{\mathcal F}(\mathbf U_e(s_r))
          \cdot\mathbf n_e
 -(\overline{\boldsymbol\sigma}_T\cdot\mathbf n_e)
        \frac{\mathbf U_a+\mathbf U_b}{2}\right].
 \label{eq:ALE_RD_Gauss_new_edge}
\end{align}
The mesh-advection term is unchanged: with the element mean mesh velocity
it is linear in $\mathbf U_e$. This statement concerns the spatial residual
of equation~\eqref{eq:ALE_RD_mean_velocity_identity}. It does not replace
the actual moving-face velocity by its element mean in the conservative
ALE balance, or remove the geometric term represented by the modified mass.

For the Gauss--Legendre rule on $[0,1]$, let $z_r$ be the roots of the
Legendre polynomial $P_q$. Then
\begin{equation}
 s_r=\frac{1+z_r}{2},\qquad
 w_r=\frac{1}{(1-z_r^2)[P_q'(z_r)]^2}.
\end{equation}
The three-point rule has
\begin{equation}
 (s_1,s_2,s_3)=\left(\frac12-\frac{\sqrt{15}}{10},\frac12,
                          \frac12+\frac{\sqrt{15}}{10}\right),\qquad
 (w_1,w_2,w_3)=\left(\frac5{18},\frac49,\frac5{18}\right).
\end{equation}
The three- and five-point rules integrate polynomials through degrees five
and nine exactly. Euler fluxes contain ratios such as $m_x^2/\rho$ and
$(E+p)m_x/\rho$, so these rules do not generally integrate the Euler flux
exactly. Nor do they raise the order of the complete $P^1$/RK2 method to
five or nine.

The implementation may add the correction
\begin{align}
 \Delta\Phi_G^T
 &=\widetilde\Phi_{G_q}^T-\widetilde\Phi_{\mathrm{nodal}}^T\nonumber\\
 &=\sum_e\ell_e\sum_rw_r\Big[
 \boldsymbol{\mathcal F}(\mathbf U_e(s_r))
 -(1-s_r)\boldsymbol{\mathcal F}(\mathbf U_a)
 -s_r\boldsymbol{\mathcal F}(\mathbf U_b)\Big]\cdot\mathbf n_e.
 \label{eq:ALE_RD_Gauss_correction}
\end{align}
It cancels between neighbouring triangles sharing the same edge states.
Its mass component vanishes because the physical mass flux is linear in
the conservative momentum. The momentum and energy corrections can still
affect every component through the subsequent coupled LDA distribution.
The number of nodal unknowns, the $P^1$ basis, the F1 temporal mass choice,
the two RK stages and the topology-repair procedure are unchanged.

For illustration, even the scalar flux $f(u)=u^2/2$ evaluated on
$u_e(s)=a+(b-a)s$ gives
\begin{equation}
 \int_0^1f(u_e(s))\,\mathrm ds-\frac{f(a)+f(b)}2
 =-\frac{(b-a)^2}{12}.
\end{equation}
A linear state therefore need not produce a linear flux. Better quadrature
does not in itself define a positive dissipation operator. Although convex
interpolation of admissible ideal-gas conservative states remains
admissible at the integration points, this is not a positivity or entropy
stability proof for the complete RD update.

\paragraph{Comparison with the legacy Roe-state split.}
Let $\mathbf K_i^{\mathrm{phys}}$ be the spatial Roe matrix before the
mean mesh-velocity shift, and let $\mathbf K_i^{\mathrm{ALE}}$ include
that shift. The older mixed-state total is
\begin{align}
 \widetilde\Phi_{\mathrm{split}}^T
 &=\sum_i\mathbf K_i^{\mathrm{ALE}}\widehat{\mathbf U}_i
 -\int_T\boldsymbol\sigma_h\cdot
       \nabla(\mathbf U_h-\widehat{\mathbf U}_h)\,\mathrm dA
 \nonumber\\
 &=\sum_i\mathbf K_i^{\mathrm{phys}}\widehat{\mathbf U}_i
 -\int_T\boldsymbol\sigma_h\cdot\nabla\mathbf U_h\,\mathrm dA.
 \label{eq:ALE_RD_Gauss_legacy_split_comparison}
\end{align}
Thus the split, nodal-flux contour and Gauss contour share the intended
conservative-$\mathbf U_h$ mesh transport but use different physical-flux
constructions. The mixed-state form remains a meaningful numerical
alternative; it is not excluded merely by writing the method in a
co-moving frame. Its split correction must be carried consistently into
that frame. A comparison must retain the same IC, mass pair, RK stages,
mesh motion and topology treatment, and distinguish physical interpolation
from coordinate conditioning.

\red{\textbf{Validation boundary.} Current local tests show reduced Sod
shock-plateau scatter and improved long-time survival of smooth KH with
Gauss contour LDA. They do not establish general positivity, nonlinear KH
convergence, or hierarchical/MPI validity. A fresh matched comparison of
the legacy split in laboratory and element frames is recorded in the
development log. Also, for singular $\mathbf S^-$, the LDA sum is
$\mathbf S^-(\mathbf S^-)^\dagger\widetilde\Phi^T$; a residual outside
its range is a separate spatial-distribution conservation problem, not a
conflict between the $P^1$ basis and quadrature.}

